\section{Fiber-product identification}
\label{sec:fiber-product}

The invariants of Section~\ref{sec:group-structure} suggest that the
native automorphism group is obtained by coupling an $S_5$ coordinate
to a dihedral coordinate of order $8$ through compatible quotient
characters.

We use the presentation
\[
D_8
=
\langle r,s
\mid
r^4=s^2=1,\ srs=r^{-1}
\rangle,
\]
so $D_8$ denotes the dihedral group of order $8$.

\subsection{Fiber products over \texorpdfstring{$C_2$}{C2}}

Let
\[
\operatorname{sgn}\colon S_5\longrightarrow C_2
\]
be the sign character, written additively.

Given a surjective character
\[
\chi\colon D_8\longrightarrow C_2,
\]
define
\[
F_\chi
=
S_5\times_{C_2}D_8
=
\left\{
(\sigma,d)\in S_5\times D_8:
\operatorname{sgn}(\sigma)=\chi(d)
\right\}.
\]

Because both characters are surjective,
\[
|F_\chi|
=
\frac{|S_5||D_8|}{|C_2|}
=
\frac{120\cdot8}{2}
=
480.
\]

The projections onto $S_5$ and $D_8$ are surjective. Their kernels have
orders $4$ and $60$, respectively:
\[
\ker(F_\chi\to S_5)
\cong
\ker\chi,
\]
and
\[
\ker(F_\chi\to D_8)
\cong
A_5.
\]

Thus the choice of the character $\chi$ determines how the local
dihedral coordinate is coupled to permutation parity.

\subsection{Characters of the dihedral factor}

The abelianization of $D_8$ is
\[
D_8^{\mathrm{ab}}\cong C_2\times C_2,
\]
so there are three nontrivial characters
\[
D_8\longrightarrow C_2.
\]

Up to automorphism of $D_8$, two kernel types are relevant.

One character has cyclic kernel
\[
\langle r\rangle\cong C_4.
\]

The other character type has Klein four kernel. We choose
\[
\chi_{\mathrm{V4}}\colon D_8\longrightarrow C_2
\]
so that
\[
\ker\chi_{\mathrm{V4}}
=
\langle r^2,s\rangle
\cong C_2\times C_2.
\]

Equivalently, one may take
\[
\chi_{\mathrm{V4}}(r)=1,
\qquad
\chi_{\mathrm{V4}}(s)=0.
\]

The alternative Klein four kernel obtained by replacing $s$ with $rs$
gives an equivalent model under an automorphism of $D_8$.

\subsection{The cyclic-kernel candidate}

First consider
\[
F_{\mathrm{C4}}
=
S_5\times_{\operatorname{sgn},\,\chi_{\mathrm{C4}}}D_8,
\]
where
\[
\ker\chi_{\mathrm{C4}}\cong C_4.
\]

This candidate has the correct order, but its internal subgroup structure
does not match the native group. In particular, the induced kernel over
the $S_5$ projection is cyclic of order $4$, while the native
classification requires the four-element kernel associated with the
matching quotient character to have Klein four structure.

The cyclic-kernel model is therefore rejected.

\subsection{The Klein-four-kernel candidate}

Now define
\[
F
=
S_5\times_{\operatorname{sgn},\,\chi_{\mathrm{V4}}}D_8.
\]

Thus
\[
F
=
\left\{
(\sigma,d)\in S_5\times D_8:
\operatorname{sgn}(\sigma)=\chi_{\mathrm{V4}}(d)
\right\}.
\]

This group has order $480$. Its center has order $2$, generated by
\[
(1,r^2).
\]

Its derived subgroup is
\[
F'
=
A_5\times\langle r^2\rangle
\cong
A_5\times C_2.
\]

Its second derived subgroup is
\[
F''\cong A_5,
\]
and its abelianization is
\[
F/F'
\cong C_2\times C_2.
\]

These agree with the independently computed native invariants.

The full element-order profile also agrees:
\[
\begin{array}{c|rrrrrrrr}
\text{order}
&1&2&3&4&5&6&10&12\\
\hline
\text{count}
&1&83&20&140&24&100&72&40.
\end{array}
\]

The three index-two subgroups of $F$ likewise reproduce the native
subgroup types, including one subgroup isomorphic to
\[
S_5\times C_2
\]
and one subgroup isomorphic to the sign pullback
\[
S_5\times_{\operatorname{sgn}}C_4.
\]

\subsection{Abstract classification}

The agreement of order, center, derived series, abelianization,
element-order profile, and index-two subgroup structure identifies the
correct abstract model. The next section strengthens this identification
by constructing an explicit multiplication-preserving bijection.

\begin{theorem}
\label{thm:fiber-product-identification}
The full automorphism group of the native graph is isomorphic to the
Klein-four-kernel fiber product
\[
\operatorname{Aut}(G60)
\cong
S_5\times_{C_2}D_8,
\]
where the defining compatibility condition is
\[
\operatorname{sgn}(\sigma)=\chi_{\mathrm{V4}}(d)
\]
and
\[
\ker\chi_{\mathrm{V4}}\cong C_2\times C_2.
\]
\end{theorem}

The notation
\[
S_5\times_{C_2}D_8
\]
will refer to this specific character choice for the remainder of the
paper.

The fiber product should not be read as a direct product. The $S_5$ and
$D_8$ coordinates are constrained by a shared parity. The dihedral
coordinate therefore records more than an independent local symmetry:
its quotient parity is locked to the sign of the $S_5$ coordinate.
